% answers_short.tex —— 2010 年·全国大纲卷Ⅱ·理科 · 简短答案速查
% 由 tools/make_short_answers.py 自 answers.tex 的 \gksolhead/\gkanswer 提取，
% 并与 2010/2010_数学真题_全国2理_答案要点.md 逐题比对后生成，勿手改（改源 MD 后重跑）。

\begin{gknotice}
\textbf{使用说明}：本节只给最终答案便于快速核对，完整推导见解析版。答案为编写者独立推导，\textbf{非官方标准答案}。
\end{gknotice}

\gkpart{选择题}
\begin{center}\begin{tabular}{ccccccccccccc}
\toprule
题号 & 1 & 2 & 3 & 4 & 5 & 6 & 7 & 8 & 9 & 10 & 11 & 12 \\
\midrule
答案 & A & D & C & C & C & B & B & B & C & A & D & B \\
\bottomrule
\end{tabular}\end{center}

\gkpart{填空题}
\begin{itemize}[leftmargin=2.6em]
  \item[\textbf{13.}] $-\dfrac43$
  \item[\textbf{14.}] $1$
  \item[\textbf{15.}] $2$
  \item[\textbf{16.}] $3$
\end{itemize}

\gkpart{解答题}
\begin{itemize}[leftmargin=2.6em]
  \item[\textbf{17.}] $AD=25$。
  \item[\textbf{18.}] (1) $\displaystyle\lim_{n\to\infty}\dfrac{a_n}{S_n}=\dfrac23$；(2) 见解析版
  \item[\textbf{19.}] (1) $DE$ 为异面直线 $AB_1$ 与 $CD$ 的公垂线；(2) 二面角 $A_1-AC_1-B_1$ 的大小为 $\arccos\dfrac1{\sqrt{15}}$
  \item[\textbf{20.}] (1) $p=0.9$；(2) $0.9891$
  \item[\textbf{21.}] (1) $e=2$；(2) 见解析版
  \item[\textbf{22.}] (1) 见解析版；(2) $0\leq a\leq\dfrac12$
\end{itemize}
